\section{Algorithm: ILP Redistricting Formulation}
\label{sec:algorithm}

\subsection{Notation}

Let $G_v = (V_v, E_v)$ be the subgraph at bisection tree node $v$, with
$n = |V_v|$ tracts and $k$ districts to produce.
Let $\mathrm{pop}(t)$ be the population of tract $t$, and
$P_{\mathrm{ideal}} = \sum_t \mathrm{pop}(t) / k$ the ideal district
population.
Let $\varepsilon_{\mathrm{bal}} \in [0,1]$ be the population balance
tolerance (default $0.5\%$, i.e., $\varepsilon_{\mathrm{bal}} = 0.005$).

For each district $d \in \{0, \ldots, k-1\}$, designate an arbitrary
root $r_d$ (we use $r_d = $ the tract with smallest global index, which
is deterministic and audit-compatible).

\subsection{Decision Variables}

\begin{itemize}
\item $x_{t,d} \in \{0,1\}$: equals 1 iff tract $t$ is assigned to
  district $d$.
  Binary.
  $n \times k$ variables total.

\item $z_{t,t'} \in \{0,1\}$ for each undirected edge $\{t,t'\} \in E_v$
  with $t < t'$: equals 1 iff the edge is cut (endpoints in different
  districts).
  Binary.
  $|E_v|$ variables total (one per undirected edge; shared across districts).

\item $f_{t,t',d} \geq 0$ for each directed edge $(t,t')$ (both orientations
  of each undirected edge) and each district $d$: flow on arc $(t,t')$ for
  district $d$'s connectivity spanning tree.
  Continuous.
  $2|E_v| \times k$ variables total.
\end{itemize}

\subsection{Objective}

\begin{equation}
\label{eq:objective}
\min \sum_{\{t,t'\} \in E_v,\; t < t'} z_{t,t'}.
\end{equation}

Each undirected edge is counted exactly once in the sum.
The objective counts the total number of cut edges in the partition.

\begin{remark}[Corrected objective]
An earlier formulation (pre-R1) introduced per-district per-edge auxiliary
variables $y_{t,t',d}$ and minimised $\sum_d \sum_{t<t'} y_{t,t',d}$,
which over-counts each cut edge $k$ times (once per district).
The corrected formulation uses a single indicator $z_{t,t'}$ per undirected
edge and counts each cut edge exactly once.
This correction is critical: the old formulation favoured high-$k$ solutions
by $k$-fold over-penalising each cut, distorting the optimisation landscape.
\end{remark}

\subsection{Constraints}

\paragraph{Coverage.}
Every tract is assigned to exactly one district:
\begin{equation}
\label{eq:coverage}
\sum_{d=0}^{k-1} x_{t,d} = 1 \qquad \forall t \in V_v.
\end{equation}

\paragraph{Population balance.}
Each district's population is within $\varepsilon_{\mathrm{bal}}$ of ideal:
\begin{equation}
\label{eq:balance}
(1 - \varepsilon_{\mathrm{bal}}) \cdot P_{\mathrm{ideal}}
\leq \sum_{t \in V_v} \mathrm{pop}(t) \cdot x_{t,d}
\leq (1 + \varepsilon_{\mathrm{bal}}) \cdot P_{\mathrm{ideal}}
\qquad \forall d.
\end{equation}

\paragraph{Cut-edge linearisation.}
$z_{t,t'}$ is forced to 1 whenever the two endpoints are in different
districts:
\begin{align}
z_{t,t'} &\geq x_{t,d} - x_{t',d}
  \qquad \forall \{t,t'\} \in E_v,\; t < t',\; \forall d, \label{eq:cut1} \\
z_{t,t'} &\geq x_{t',d} - x_{t,d}
  \qquad \forall \{t,t'\} \in E_v,\; t < t',\; \forall d. \label{eq:cut2}
\end{align}
Together \eqref{eq:cut1}--\eqref{eq:cut2} enforce $z_{t,t'} \geq |x_{t,d} - x_{t',d}|$
for all $d$; since $z_{t,t'}$ is minimised in the objective, it takes
value 1 if and only if the edge is cut by at least one district boundary.

\paragraph{Contiguity (MTZ flow).}
For each district $d$, the spanning-tree flow constraints enforce that
the tracts assigned to $d$ form a connected subgraph.
Root $r_d$ supplies $n_d - 1$ units of flow; each non-root tract in $d$
consumes one unit:
\begin{align}
f_{t,t',d} &\geq 0
  \qquad \forall (t,t') \in E_v^{\mathrm{dir}},\; \forall d,
  \label{eq:flow-nonneg} \\
f_{t,t',d} &\leq (n-1) \cdot \min(x_{t,d},\, x_{t',d})
  \qquad \forall (t,t') \in E_v^{\mathrm{dir}},\; \forall d,
  \label{eq:flow-cap} \\
\sum_{t': (t,t') \in E_v^{\mathrm{dir}}} f_{t,t',d}
  - \sum_{t': (t',t) \in E_v^{\mathrm{dir}}} f_{t',t,d}
  &= x_{t,d}
  \qquad \forall t \neq r_d,\; \forall d,
  \label{eq:flow-balance} \\
\sum_{t': (r_d,t') \in E_v^{\mathrm{dir}}} f_{r_d,t',d}
  &= \sum_{t \in V_v} x_{t,d} - 1
  \qquad \forall d.
  \label{eq:flow-root}
\end{align}
Here $E_v^{\mathrm{dir}}$ contains both orientations of each undirected
edge in $E_v$.

\begin{remark}[Both-endpoint capacity bound]
Constraint~\eqref{eq:flow-cap} uses $\min(x_{t,d}, x_{t',d})$ rather
than the standard $x_{t,d}$ bound.
This ensures that flow on arc $(t, t')$ is zero whenever \emph{either}
endpoint is outside district $d$: if $x_{t',d} = 0$, the flow cap
drops to zero regardless of $x_{t,d}$.
The $\min$ is linearised as two linear constraints:
$f_{t,t',d} \leq (n-1) \cdot x_{t,d}$ and
$f_{t,t',d} \leq (n-1) \cdot x_{t',d}$.
This tightens the LP relaxation relative to standard MTZ and reduces
branch-and-bound tree size~\citep{gurnee2021}.
\end{remark}

\subsection{Variable Counts and Problem Size}

For reference, Table~\ref{tab:var-counts} gives variable and constraint
counts for representative instances.

\begin{table}[t]
\centering
\caption{%
  Variable and constraint counts for representative instances.
  $|E|$ is the number of undirected edges in the tract adjacency graph.
  Typical adjacency degree is 5--8 per tract for census-tract graphs.
}
\label{tab:var-counts}
\renewcommand{\arraystretch}{1.3}
\begin{tabular}{@{}lrrrrrrr@{}}
\toprule
Instance & $n$ & $k$ & $|E|$ & Bin.\ $x$ & Bin.\ $z$ & Cont.\ $f$ & Constraints \\
\midrule
4-node path ($k=2$)   & 4   & 2 & 3  & 8   & 3   & 12  & $\approx 50$ \\
VT 2020 ($k=1$)        & 181 & 1 & 894 & 181 & 894 & 0   & $\approx 1{,}300$ \\
NH 2020 ($k=2$)        & 260 & 2 & 1{,}287 & 520 & 1{,}287 & 5{,}148 & $\approx 9{,}400$ \\
RI 2020 ($k=2$)        & 244 & 2 & 1{,}202 & 488 & 1{,}202 & 4{,}808 & $\approx 8{,}800$ \\
Bisect node ($n=400$, $k=2$) & 400 & 2 & 1{,}980 & 800 & 1{,}980 & 7{,}920 & $\approx 14{,}500$ \\
\bottomrule
\end{tabular}
\end{table}

\begin{remark}[$k=1$ degenerate case]
For $k=1$ (single-district states), there is exactly one valid assignment:
every tract goes to district 0.
The ILP reduces to a feasibility check with trivially zero edge cut.
The solver returns \texttt{Optimal} immediately with EC$= 0$ and solve
time under 1\,s regardless of state size.
This degenerate case is detected and fast-pathed before solver invocation.
\end{remark}

\subsection{Formal Pseudocode}

Algorithm~\ref{alg:ilp} gives the top-level \ilp{} redistricting procedure.

\begin{algorithm}[t]
\caption{\ilp{}-\textsc{Bisect}($G_v$, $k$, $\mathbf{p}$,
  $\varepsilon_{\mathrm{bal}}$, $\tau$, $\delta$, $n_{\max}$)}
\label{alg:ilp}
\begin{algorithmic}[1]
\State $n \gets |V_v|$
\If{$k = 1$}
  \textbf{return} $(\text{all tracts} \to 0,\; \texttt{Optimal},\; \text{EC}=0,\; \text{gap}=0.0)$
\EndIf
\If{$n > n_{\max}$}
  \State $\mathbf{warn}$(\texttt{"ILP fallback: $n > n_{\max}$, using METIS"})
  \State \textbf{return} \metis{}-\textsc{Bisect}$(G_v, k, \mathbf{p}, \varepsilon_{\mathrm{bal}})$
  \Comment{size guard}
\EndIf
\State Build LP/MPS file $F$ with constraints
  \eqref{eq:coverage}--\eqref{eq:flow-root}
  \Comment{write to \texttt{runs/\ldots/ilp/\{node\}.mps}}
\State $(x^*, z^*, \text{status}, \text{gap}) \gets \textsc{Solve}(F, \tau, \delta)$
  \Comment{$\tau$: time limit, $\delta$: gap threshold}
\If{status $\in$ \{\texttt{Optimal}, \texttt{GapReached}\}}
  \State $\mathrm{assign}(t) \gets \arg\max_d x^*_{t,d}$ for each $t$
  \State \textbf{return} $(\mathrm{assign},\; \text{status},\; \sum_{t<t'} z^*_{t,t'},\; \text{gap})$
\ElsIf{status $=$ \texttt{Timeout}}
  \State \textbf{return} best incumbent with recorded gap $> \delta$
\Else
  \textbf{raise} \textsc{Infeasible}
\EndIf
\end{algorithmic}
\end{algorithm}

\subsection{Solver Selection}

The solver interface is abstracted behind a \texttt{Solver} trait so that
backends can be swapped without touching the formulation code.
Phase~1 uses a two-tier priority order:

\begin{enumerate}
\item \textbf{\glpk{} via \texttt{good\_lp}} --- pure Rust $+$ C binding,
  no subprocess, bundled with the crate.
  Default for $n \leq 200$.
  Typical solve time for NH: 47\,s.

\item \textbf{\highs{} subprocess} --- invokes the \texttt{highs} binary
  from \texttt{PATH} if available.
  For $200 < n \leq 500$.
  \highs{} is a state-of-the-art open-source MIP solver with significantly
  faster branch-and-bound than \glpk{} on larger instances.
\end{enumerate}

If neither solver is available and $n \leq n_{\max}$, the MPS file is
written and the CLI returns an error directing the user to install
\texttt{highs}.

\subsection{Size Guard and Fallback}

If the subgraph at a bisection node has $n > \texttt{max\_tracts}$
(default 500), \ilp{} is not invoked and \metis{} runs instead:

\begin{verbatim}
WARNING: ILP solver skipped for node at depth 2 (847 tracts > 500 limit).
         Falling back to METIS.
         Increase --ilp-max-tracts or use --structure metis to suppress.
\end{verbatim}

The fallback is recorded in the audit JSON as
\texttt{"solver\_status": "fallback\_metis"} so that the audit chain
accurately reflects which portions of the bisection tree used exact
optimisation and which used heuristic fallback.

\subsection{MPS File Retention}

Before invoking the solver, the formulation is serialised to
\texttt{runs/\{label\}/\{year\}/ilp/\{node\_id\}.mps} in the standard
Mathematical Programming System format.
This file is retained regardless of solver outcome.
An external auditor can reload it into any LP/MIP solver and verify that
the committed plan achieves the reported optimal objective value.
This is the practical manifestation of the ``no valid plan has lower
edge cuts'' certificate: the certificate is the MPS file plus the
solver's return status.
